%% kapa: exam=1 round=36 year=2025 period=1 form=A subject=civil_law subject_ko=민법 seq=19 number=19
%% source: Q-Net 감정평가사 기출문제 (cst00309) artlSeq=5231525 file=raw/1008/5231525/2025년 제36회 감정평가사 1차시험 1교시 문제지 원본.pdf
%% conversion: kordoc 4.24.51 -> tools/gen1.py (auto); review=proofread-sonnet; flags=none
\begin{question}{19}
사회질서에 반하는 법률행위로서 무효인 것은? (다툼이 있으면 판례에 따름)
\begin{choices}
  \item 양도소득세의 일부를 회피할 목적으로 매매계약서에 실거래가액보다 낮은 금액을 매매대금으로 기재한 매매계약
  \item 강제집행을 면할 목적으로 부동산에 허위의 근저당권설정등기를 경료하는 행위
  \item 반사회적 행위에 의하여 조성된 비자금을 소극적으로 은닉하기 위하여 임치한 행위
  \item 전통사찰의 주지직을 거액의 금품을 대가로 양도하기로 하는 약정이 있음을 알고도 이를 묵인한 상태에서 한 종교법인의 주지임명행위
  \item 보험계약자가 다수의 보험계약을 통하여 보험금을 부정취득할 목적으로 체결한 보험계약
\end{choices}
\end{question}
\answer{5}
